\chapter{Preface}

Garry Kasparov wrote a book about how life imitates chess~\cite{Kasparov:2007-life}. While studying
endgames for this book, I found something similar. Three habits from the endgame have helped me
far beyond the board.

\begin{itemize}
\item Prophylaxis. Before making your own plans, ask what could go wrong, and remove the potential
damage first. In life this means, for example, not using leverage: a risk you cannot afford is not
worth any gain.

\item The process of elimination. When the best move is not clear, rule out the moves that lose.
What remains is usually good enough. Many decisions in life become easier the same way.

\item Typical plans. Every endgame has its typical resources, the defensive setups and the attacking
plans listed at the start of each chapter. Once you know them, you recognise the situation and know
what to aim for.
\end{itemize}

When we face a similar situation in life, we can orient ourselves in the same way: first learn the
situation, then check whether one of the typical plans works for it.

I hope this book helps you, as it has helped me, to play better chess, and perhaps to navigate life
a little better.

\begin{flushright}
Ming Lu
\end{flushright}
